\chapter{Methodology}
\label{ch:methodology}

\section{Overview}
\label{sec:methodology-overview}

Our approach addresses key-value pair extraction as a two-stage discriminative
problem: (i)~\emph{entity classification}, which assigns each text token a role
(\textsc{key}, \textsc{value}, or \textsc{other}); and (ii)~\emph{linking},
which pairs each predicted key span with its corresponding value span.
Both tasks are solved jointly by a single LayoutLMv3-based model augmented
with a biaffine linker head, trained end-to-end on the KVP10k dataset.

We investigate two linking granularities: \textbf{V1 (token-level)}, where
every key token is scored against every value token, and \textbf{V2
(span-level)}, where contiguous same-label tokens are first grouped into
spans before pairwise scoring. The span-level approach is motivated by
SPADE~\cite{spade} and BROS~\cite{bros}, which demonstrate that entity
linking operates naturally at the segment level (Section~\ref{subsec:entity_linking}).

Figure~\ref{fig:architecture-overview} provides a schematic overview.
The encoder receives token sequences together with 2D bounding-box coordinates
and document image patches. Its output representations are consumed in parallel
by the entity-classification head and the biaffine linker head.

\begin{figure}[ht]
  \centering
  \begin{tikzpicture}[
    block/.style={rectangle, draw, rounded corners, minimum height=1cm, minimum width=2.8cm, align=center, font=\small},
    io/.style={rectangle, draw, minimum height=0.8cm, minimum width=2.2cm, align=center, font=\small, fill=gray!10},
    arrow/.style={-{Stealth[length=3mm]}, thick},
    node distance=0.8cm and 1.2cm
  ]
    % Input
    \node[io] (input) {OCR Words +\\Bounding Boxes +\\Document Image};

    % Encoder
    \node[block, below=of input, fill=blue!10] (encoder) {LayoutLMv3\\Encoder};

    % Two heads
    \node[block, below left=1.2cm and 1.5cm of encoder, fill=orange!15] (entity) {Entity Classifier\\(Key / Value / Other)};
    \node[block, below right=1.2cm and 1.5cm of encoder, fill=red!12] (linker) {Biaffine Linker\\(Key--Value Scoring)};

    % Span grouper (V2 only)
    \node[block, below=0.8cm of entity, fill=green!12] (spans) {Span Grouper\\(V2 only)};

    % Output
    \node[io, below=2.5cm of encoder] (output) {KVP Pairs\\(key $\rightarrow$ value)};

    % Arrows
    \draw[arrow] (input) -- (encoder);
    \draw[arrow] (encoder) -| (entity);
    \draw[arrow] (encoder) -| (linker);
    \draw[arrow] (entity) -- (spans);
    \draw[arrow] (spans) -| (linker);
    \draw[arrow] (linker) |- (output);

    % Labels
    \node[font=\scriptsize\itshape, right=0.1cm of spans] {contiguous merge};
  \end{tikzpicture}
  \caption{Overview of the proposed LayoutLMv3 + Biaffine Linker architecture.
    Token representations produced by the encoder are fed to two heads:
    an entity classifier (token-level labelling) and a biaffine linker
    (pairwise key--value scoring). In V2, a span grouper merges contiguous
    same-label tokens into spans before linking.}
  \label{fig:architecture-overview}
\end{figure}

% ---------------------------------------------------------------------------
\section{LayoutLMv3 Encoder}
\label{sec:layoutlmv3-encoder}

\subsection{Input Representation}
\label{sec:input-representation}

LayoutLMv3~\cite{huang2022layoutlmv3} accepts three parallel input streams that
are fused inside the transformer backbone.

\paragraph{Text tokens.}
The OCR word sequence is tokenised with the LayoutLMv3 tokenizer (a byte-pair
encoding vocabulary of 50{,}265 entries). A \texttt{[CLS]} token is prepended
and a \texttt{[SEP]} token appended. For multi-word entities, the bounding box
of the first sub-word token is used for the whole entity.

\paragraph{Layout embeddings.}
Each token is associated with its bounding box on a normalised $1000 \times 1000$
coordinate grid. Following the LayoutLMv3 convention, six scalar values are
derived per token:
\[
  \mathbf{b}_i = (x_{\min},\, y_{\min},\, x_{\max},\, y_{\max},\, w,\, h)
\]
where $w = x_{\max} - x_{\min}$ and $h = y_{\max} - y_{\min}$. Each scalar
is embedded independently with a learnable embedding table and the six
embeddings are summed with the text token embedding before the first
transformer layer.

\paragraph{Image patch embeddings.}
The document image is resized to $224 \times 224$ pixels and split into
$16 \times 16 = 196$ non-overlapping patches of size $14 \times 14$.
Each patch is projected to the model hidden dimension $d = 768$ via a
learnable linear layer. A 2D sine-cosine positional encoding is added to
preserve patch location. The resulting 196 patch tokens are concatenated with
the text tokens to form the full input sequence.

\subsection{Self-Attention Mechanism}
\label{sec:self-attention}

LayoutLMv3 uses a standard multi-head self-attention with $H = 12$ heads and
hidden dimension $d = 768$. The spatial-aware bias introduced in
LayoutLMv2~\cite{xu2021layoutlmv2} is \emph{not} re-applied; instead, layout
information is injected through the embedding summing described above, allowing
the standard dot-product attention to propagate spatial signals across layers.
Cross-modal attention between text tokens and image patch tokens is unrestricted,
enabling the model to ground textual entities in specific image regions.

\subsection{Pre-training Objectives}
\label{sec:pretraining-objectives}

The LayoutLMv3 checkpoint used in this work
(\texttt{microsoft/layoutlmv3-base}) was pre-trained on IIT-CDIP with three
objectives~\cite{huang2022layoutlmv3}:

\begin{enumerate}
  \item \textbf{Masked Language Modelling (MLM):} 15\% of text tokens are masked
  and the model predicts the missing tokens, encouraging semantic text understanding.
  \item \textbf{Masked Image Modelling (MIM):} 40\% of image patches are masked
  and the model reconstructs the discrete token index of each patch (using a
  pre-trained image tokenizer), encouraging visual-spatial grounding.
  \item \textbf{Word-Patch Alignment (WPA):} The model predicts whether a text
  token and an image patch originate from the same document region, creating
  explicit cross-modal alignment supervision.
\end{enumerate}

These pre-training objectives give the encoder strong initialisation for both
textual and layout-spatial understanding, reducing the amount of task-specific
fine-tuning data required.

% ---------------------------------------------------------------------------
\section{Task Head: Entity Classifier}
\label{sec:entity-classifier}

Rather than a sequence-to-sequence decoder, we frame extraction as token
classification. The entity head maps each text token representation to one of
three classes: \texttt{O} (other), \texttt{KEY}, \texttt{VALUE}.

\paragraph{Architecture.}
Let $\mathbf{H} \in \mathbb{R}^{N \times d}$ denote the encoder output for
$N$ text tokens. A two-layer MLP with GELU activations projects each token to
logits:
\[
  \hat{y}_i = \operatorname{softmax}\!\bigl(\mathbf{W}_2\,
    \operatorname{GELU}(\mathbf{W}_1\,\mathbf{h}_i + \mathbf{b}_1) +
    \mathbf{b}_2\bigr)
\]
where $\mathbf{W}_1 \in \mathbb{R}^{d \times d}$,
$\mathbf{W}_2 \in \mathbb{R}^{d \times 3}$. Dropout with $p = 0.1$ is applied
after the first projection.

\paragraph{Why token classification instead of generation?}
Generative baselines (Stage~3 Mistral) demonstrated strong text recall but
poor spatial grounding and systematic over-prediction on sparse layouts
(Section~\ref{sec:per-cluster-analysis}). Token classification is grounded by
construction — predictions are tied to OCR tokens with known bounding boxes —
eliminating the spatial anchoring problem entirely.

% ---------------------------------------------------------------------------
\section{Biaffine Linker (V1: Token-Level)}
\label{sec:biaffine-linker}

\subsection{Scoring Function}
\label{sec:kv-attention}

After entity labels are predicted, the linker determines which \textsc{key}
tokens should be paired with which \textsc{value} tokens. The scoring function
combines a semantic signal with a spatial signal.

\paragraph{Semantic scorer.}
The semantic score uses a scaled dot-product between projected key and value
representations:
\[
  s_{\text{sem}}(i,j) = \frac{(\mathbf{W}_k\,\mathbf{h}_i)^\top\,
  (\mathbf{W}_v\,\mathbf{h}_j)}{\sqrt{d}}
\]
where $\mathbf{h}_i, \mathbf{h}_j \in \mathbb{R}^{d}$ are encoder hidden
states ($d = 768$), and $\mathbf{W}_k, \mathbf{W}_v \in \mathbb{R}^{d \times d}$
are learnable projections. This formulation follows the attention mechanism
of Vaswani et al.~\cite{vaswani2017attention} and is numerically stable
because no learned weight matrix couples the key and value streams in the
scoring step.

\paragraph{Design evolution.}
An earlier implementation used a full biaffine formulation
($s = \mathbf{k}^\top \mathbf{W}_{\text{bil}} \mathbf{v}$, following
Dozat and Manning~\cite{dozat2017biaffine}), but this caused
\texttt{link\_loss~=~NaN} at approximately step~1400 due to exploding
gradients through the $768 \times 768$ bilinear weight matrix. The scaled
dot-product formulation eliminates this failure mode.

\subsection{Spatial Encoding}
\label{sec:spatial-encoding}

Eight geometric features are computed per (key, value) token pair: horizontal
and vertical centre displacement $(\Delta x, \Delta y)$, Euclidean distance,
angle $\theta = \text{atan2}(\Delta y, \Delta x)$, horizontal and vertical
bounding-box overlap ratios, relative area ratio, and relative aspect-ratio
ratio. These are encoded by a two-layer MLP with ReLU activations into a
32-dimensional vector, which is concatenated with the semantic scalar before
a final MLP produces the link score $s_{i,j}$.

\subsection{Class Imbalance Problem}
\label{sec:v1-imbalance}

In V1, every key token is scored against every value token, producing a dense
$n_k \times n_v$ matrix. On a typical KVP10k page with $\sim$100 key tokens
and $\sim$100 value tokens, this yields $\sim$10{,}000 candidate pairs of
which only $\sim$20 are positive (true links), giving a negative-to-positive
ratio of approximately \textbf{450:1}. Despite applying
\texttt{pos\_weight}~$= \min(n_{\text{neg}}/n_{\text{pos}},\, 50)$ to the
BCE loss, the extreme imbalance prevents the linker from learning meaningful
associations, resulting in link F1~$\approx 0.008$ (see
Section~\ref{sec:v1-results}).

This observation motivates V2: by grouping tokens into spans before linking,
the candidate matrix shrinks from $\sim$10{,}000 to $\sim$225 entries with
a manageable $\sim$5:1 imbalance ratio.

% ---------------------------------------------------------------------------
\section{Span-Level Linker (V2)}
\label{sec:v2-span-linking}

V2 addresses the granularity mismatch identified in V1 by performing linking
at the \emph{span} level rather than the token level. This design is directly
motivated by SPADE~\cite{spade}, which operates on text segments, and
BROS~\cite{bros}, which links entity-initial tokens to reduce the prediction
space.

\subsection{Span Grouping}
\label{sec:span-grouping}

After entity classification, contiguous tokens sharing the same non-\texttt{O}
label are merged into spans. For example, if tokens 5--7 are labelled
\texttt{KEY} and tokens 8--12 are labelled \texttt{VALUE}, they form one key
span and one value span respectively. This grouping is performed
deterministically at both training and inference time.

Figure~\ref{fig:span-grouping} illustrates the span grouping process.

\begin{figure}[ht]
  \centering
  \begin{tikzpicture}[
    tok/.style={rectangle, draw, minimum height=0.7cm, minimum width=1.2cm,
                align=center, font=\scriptsize},
    span/.style={rectangle, draw, rounded corners, minimum height=0.8cm,
                 align=center, font=\scriptsize\bfseries},
    arrow/.style={-{Stealth[length=2mm]}, thick, gray},
    node distance=0.05cm
  ]
    % Token row
    \node[tok, fill=gray!15] (t1) {Invoice};
    \node[tok, fill=red!20, right=of t1] (t2) {Date};
    \node[tok, fill=red!20, right=of t2] (t3) {:};
    \node[tok, fill=blue!20, right=of t3] (t4) {01};
    \node[tok, fill=blue!20, right=of t4] (t5) {/15/};
    \node[tok, fill=blue!20, right=of t5] (t6) {2024};
    \node[tok, fill=red!20, right=of t6] (t7) {Total};
    \node[tok, fill=red!20, right=of t7] (t8) {Amt};
    \node[tok, fill=blue!20, right=of t8] (t9) {\$1500};

    % Labels
    \node[font=\tiny, above=0.15cm of t1] {O};
    \node[font=\tiny, above=0.15cm of t2] {KEY};
    \node[font=\tiny, above=0.15cm of t3] {KEY};
    \node[font=\tiny, above=0.15cm of t4] {VAL};
    \node[font=\tiny, above=0.15cm of t5] {VAL};
    \node[font=\tiny, above=0.15cm of t6] {VAL};
    \node[font=\tiny, above=0.15cm of t7] {KEY};
    \node[font=\tiny, above=0.15cm of t8] {KEY};
    \node[font=\tiny, above=0.15cm of t9] {VAL};

    % Span row
    \node[span, fill=red!30, below=1.2cm of $(t2)!0.5!(t3)$] (s1) {Key 1:\\``Date :''};
    \node[span, fill=blue!30, below=1.2cm of t5] (s2) {Val 1:\\``01/15/2024''};
    \node[span, fill=red!30, below=1.2cm of $(t7)!0.5!(t8)$] (s3) {Key 2:\\``Total Amt''};
    \node[span, fill=blue!30, below=1.2cm of t9] (s4) {Val 2:\\``\$1500''};

    % Arrows
    \draw[arrow] (t2.south) -- (s1.north);
    \draw[arrow] (t3.south) -- (s1.north);
    \draw[arrow] (t4.south) -- (s2.north);
    \draw[arrow] (t5.south) -- (s2.north);
    \draw[arrow] (t6.south) -- (s2.north);
    \draw[arrow] (t7.south) -- (s3.north);
    \draw[arrow] (t8.south) -- (s3.north);
    \draw[arrow] (t9.south) -- (s4.north);
  \end{tikzpicture}
  \caption{Span grouping: contiguous tokens with the same entity label are
    merged into spans. The span representation is the mean pool of constituent
    token hidden states; the span bounding box is the union of constituent
    token boxes.}
  \label{fig:span-grouping}
\end{figure}

\subsection{Span Representation}
\label{sec:span-representation}

Each span's representation is formed by mean-pooling the encoder hidden states
of its constituent tokens:
\[
  \mathbf{s}_k = \frac{1}{|T_k|} \sum_{i \in T_k} \mathbf{h}_i
\]
where $T_k$ is the set of token indices in span $k$. The span bounding box
is the tight union box over all constituent token boxes. The same spatial
augmentation (8 geometric features, Section~\ref{sec:spatial-encoding}) is
applied at the span level.

\subsection{Span-Level Scoring}
\label{sec:span-scoring}

The span-level biaffine scorer operates identically to V1
(Section~\ref{sec:kv-attention}) but over span representations instead of
token representations. For $K$ key spans and $V$ value spans, the linker
produces a $K \times V$ score matrix. On typical KVP10k pages, $K \approx 15$
and $V \approx 15$, yielding $\sim$225 candidate pairs with
$\sim$15 positives --- a \textbf{$\sim$5:1} neg:pos ratio compared to V1's
$\sim$450:1.

\subsection{Training Objective}
\label{sec:v2-training}

V2 uses the same joint loss as V1:
\[
  \mathcal{L} = \mathcal{L}_{\text{entity}} + \lambda\,\mathcal{L}_{\text{link}}
  \tag{3.1}
\]
The entity loss is unchanged. The link loss is binary cross-entropy over the
$K \times V$ span-level score matrix. Because the class imbalance is
manageable at this granularity, \texttt{pos\_weight} is computed dynamically
per batch as $n_{\text{neg}}/n_{\text{pos}}$ without the cap of~50 required
in V1.

% ---------------------------------------------------------------------------
\section{Training Procedure}
\label{sec:training-procedure}

\subsection{Loss Function}
\label{sec:loss-function}

The model is trained with a joint loss combining entity classification and
linking (Equation~3.1, repeated here for convenience):
\[
  \mathcal{L} = \mathcal{L}_{\text{entity}} + \lambda\,\mathcal{L}_{\text{link}}
\]
where:
\begin{itemize}
  \item $\mathcal{L}_{\text{entity}}$ is the cross-entropy loss over
  \{\texttt{O, KEY, VALUE}\} for each text token, with inverse-frequency
  class weights to address the dominance of the \texttt{O} class.
  \item $\mathcal{L}_{\text{link}}$ is the binary cross-entropy loss over
  predicted link scores, with \texttt{pos\_weight} to address class imbalance
  (capped at 50 for V1; uncapped for V2).
  \item $\lambda$ is a weighting hyperparameter swept over
  $\{0.5, 1.0, 2.0\}$ during development (see Section~\ref{sec:stage4-hyperparams}).
\end{itemize}

\paragraph{Training stability.}
Achieving stable Stage~4b training required addressing two sources of
numerical instability: (i)~the switch from a bilinear weight matrix to
scaled dot-product scoring (Section~\ref{sec:kv-attention}), and
(ii)~a valid-label mask applied before the BCE call to exclude positions
with uninitialised link labels. After these fixes, training proceeds
without \texttt{NaN} across all runs. Full details of the diagnostic
process are documented in Section~\ref{sec:linker-impl}.

\subsection{Optimisation}
\label{sec:optimisation}

We use AdamW~\cite{loshchilov2019decoupled} with weight decay $0.01$ and a
linear warmup of 500 steps followed by linear decay to zero. The effective
batch size is 32, achieved through gradient accumulation over 8 steps with a
per-device batch size of 4. Training runs for a maximum of 20 epochs with
early stopping (patience~$= 3$) on the validation entity F1 score.

All Stage~4 training uses fp32 precision. An earlier attempt with bfloat16
caused \texttt{loss~=~NaN} due to underflow in the spatial feature encoder;
reverting to fp32 resolved the issue with negligible memory impact given
the compact model size (125M parameters).

\subsection{Two-Phase Training}
\label{sec:two-phase}

Training proceeds in two phases:
\begin{enumerate}
  \item \textbf{Stage~4a (entity pre-training):} The encoder and entity
  classification head are trained with $\lambda = 0$ (linker disabled).
  This establishes strong entity representations before the linking
  objective is introduced.
  \item \textbf{Stage~4b (joint fine-tuning):} The model resumes from the
  Stage~4a checkpoint with the linker activated ($\lambda > 0$). The encoder
  learning rate is reduced by a factor of 10 to preserve entity classification
  capability.
\end{enumerate}

% ---------------------------------------------------------------------------
\section{Inference}
\label{sec:inference}

\subsection{Entity Extraction}
\label{sec:greedy-decoding}

At inference time, the entity head produces per-token logits which are
argmax-decoded to \{\texttt{O, KEY, VALUE}\} labels. Consecutive tokens with
the same non-\texttt{O} label are merged into spans. The bounding box of a
span is the tight union box over all constituent token boxes.

\subsection{Linking (V1)}
\label{sec:v1-linking-inference}

In V1, the linker scores all key--value token pairs. A sigmoid threshold of
0.5 determines valid links; each key is matched to at most one value (the
highest-scoring valid pair).

\subsection{Linking (V2)}
\label{sec:v2-linking-inference}

In V2, spans are first formed by the grouping step
(Section~\ref{sec:span-grouping}), then span representations are computed via
mean-pooling. The span-level biaffine scorer produces a $K \times V$ score
matrix and the highest-scoring match for each key span is selected.

\subsection{Post-processing}
\label{sec:post-processing}

Unmatched keys are emitted as \emph{unvalued-key} annotations and unmatched
values as \emph{unkeyed-value} annotations, matching the KVP10k annotation
convention. No beam search or constrained decoding is applied: the
discriminative formulation avoids the need for autoregressive search entirely.
Bounding boxes are inherited directly from the OCR input, eliminating the
hallucinated bounding boxes that arise in generative approaches.
